% Do not change document class, margins, fonts, etc.
\documentclass[a4paper,oneside,bibliography=totoc]{scrbook}


\usepackage[utf8]{inputenc}
\usepackage{graphicx}
\usepackage{latexsym}
\usepackage{amsmath}
\usepackage{amssymb}
\usepackage{tabularx}
\usepackage{booktabs}
\usepackage{algorithm}
\counterwithin{algorithm}{chapter}
\usepackage{algorithmic}
\usepackage{csquotes}
\renewcommand{\algorithmiccomment}[1]{\hfill\textit{// #1}}
\usepackage[usenames,dvipsnames,table]{xcolor}
\usepackage{array}
\usepackage{longtable}
\usepackage{xurl}
\usepackage[colorlinks,citecolor=Green]{hyperref}

\newcolumntype{L}[1]{>{\raggedright\arraybackslash\hspace{0pt}}p{#1}}
\newcolumntype{C}[1]{>{\centering\arraybackslash}p{#1}}
\definecolor{TableHeader}{RGB}{0,51,102}
\definecolor{LightTable}{RGB}{235,241,248}
\definecolor{Missing}{RGB}{192,0,0}
\newcommand{\tablegroup}[1]{\textbf{#1}\ }
\setlength{\emergencystretch}{3em}
\newcommand{\miss}[1]{#1\textsuperscript{\textcolor{Missing}{**}}}
\usepackage{lipsum}


\usepackage{natbib}
\bibliographystyle{chicagoa}
\setcitestyle{authoryear,round,semicolon,aysep={},yysep={,}} \let\cite\citep


\newtheorem{definition}{Definition} \newtheorem{proposition}{Proposition}

\begin{document}


\makeatletter
\renewcommand\thesection{\arabic{section}}
\renewcommand\thesubsection{\arabic{section}.\arabic{subsection}}
\renewcommand\thesubsubsection{\arabic{section}.\arabic{subsection}.\arabic{subsubsection}}
\makeatother

\subject{Web Data Integration \\ Final Report}
\title{Pharmaceutical Drug Integration: Linking US, EU and Canadian Approval Data}
\author{
  Gabriel Himmelein (1649181)\\
  Özlem Ölçer (2265825)\\
  Anatole Lobenko (2187322)\\
  Hai Nguyen (2267317)\\
  Thu Huyen Trang Trinh (2111810)
} \date{November 29, 2026}
\publishers{{\small Submitted to}\\
  Data and Web Science Group\\
  Dr. Ralph Peeters\\
  Prof. Dr. Christian Bizer\\
  University of Mannheim\\}
\maketitle

% ============================================================
\section{Use Case}
% ============================================================

This project builds a unified view of approved pharmaceutical drugs across the United States, European Union, and Canada by integrating three official regulatory datasets: Drugs@FDA (US), EMA Medicines (EU), and the Health Canada Drug Product Database (DPD). The central entity is a medicine identified by its brand name and active substance(s); each integrated record includes licence holder, regulatory status, product class, ATC code, dosage form, route, strengths, substance count, regional dates, and generic/biosimilar/orphan flags.

% ============================================================
\section{Requirements Suitability}
% ============================================================

\subsection{Schema and basic profile}

\begin{table}[htbp]
\centering
\caption{Overview of the input datasets}
\label{tab:datasets}
\scriptsize
\setlength{\tabcolsep}{4pt}
\begin{tabular}{|l|l|l|r|r|r|}
\hline
\textbf{Dataset} & \textbf{Source} & \textbf{Format} & \textbf{Rows} & \textbf{Entities} & \textbf{Attr.} \\
\hline
Drugs@FDA  & U.S. FDA    & TXT/TSV  & 51,839 &  8,175 & 56 \\
EMA Med.   & EMA         & XLSX/CSV &  2,746 &  2,648 & 39 \\
DPD        & Health Can. & TXT/CSV  & 42,042 & 41,636 & 87 \\
\hline
\end{tabular}
\end{table}

Table~\ref{tab:datasets} summarises the three input datasets. Drugs@FDA contains 51,839 product rows forming 8,175 entities across 29,377 applications (78.3\% ANDAs). EMA Medicines contains 2,746 rows consolidated into 2,648 entities (including 28 records recovered via INN fallback), predominantly human medicines (85.6\%). The Health Canada DPD contains 58,137 product records consolidated into 41,636 entities. High missingness in all three datasets is largely structural: optional fields, category-specific attributes, or regulatory-event dates that apply only to specific statuses.

\subsection{Integrated schema and overlap with input schemata}

The integrated schema contains 16 attributes (Table~\ref{tab:schema}); eleven occur in at
least two sources. \texttt{activeSubstances} and \texttt{strengths} are list-valued, while
the remaining 14 attributes are single-valued (including the numeric attribute
\texttt{activeSubstanceCount}). Strengths are represented as numeric value--unit pairs.

\begin{table}[htbp]
\centering
\caption{Attribute Intersection with Integrated Schema}
\label{tab:schema}
\scriptsize
\begin{tabular}{|l|l|>{\raggedright\arraybackslash}p{5.5cm}|}
\hline
\textbf{Attribute name} & \textbf{Attribute type} & \textbf{Datasets in which the attribute is found} \\
\hline
\texttt{brandName} & string & Drugs@FDA, EMA Medicines, DPD \\
\texttt{activeSubstances} & list (string) & Drugs@FDA, EMA Medicines, DPD \\
\texttt{activeSubstanceCount} & numeric (integer) & Drugs@FDA (derived), EMA Medicines (derived), DPD \\
\texttt{licenceHolder} & string & Drugs@FDA, EMA Medicines, DPD \\
\texttt{regulatoryStatus} & string (categorical) & Drugs@FDA, EMA Medicines, DPD \\
\texttt{productClass} & string (categorical) & Drugs@FDA (derived), EMA Medicines, DPD \\
\texttt{atcCode} & string & EMA Medicines, DPD \\
\texttt{generic} & boolean & Drugs@FDA (derived), EMA Medicines \\
\texttt{biosimilar} & boolean & EMA Medicines \\
\texttt{orphanMedicine} & boolean & EMA Medicines \\
\texttt{dosageForm} & string & Drugs@FDA, DPD \\
\texttt{route} & string & Drugs@FDA, DPD \\
\texttt{strengths} & list (numeric value, unit) & Drugs@FDA, DPD \\
\texttt{usApprovalDate} & date & Drugs@FDA \\
\texttt{euAuthorisationDate} & date & EMA Medicines \\
\texttt{caStatusDate} & date & DPD \\
\hline
\end{tabular}
\end{table}

% ============================================================
\subsection{Entity Identification and Count Requirement}
\label{subsec:entity_id}
% ============================================================

A medicine is identified by its \texttt{brandName} and \texttt{activeSubstances}. Contrary
to initial assumptions, \texttt{licenceHolder} cannot be used as an identifying attribute
because the same medicine is frequently marketed by different regional subsidiaries or
unrelated licensing partners across jurisdictions (for example, \textit{Verkazia} is held
by \texttt{HARROW EYE} in the US, \texttt{Santen Oy} in the EU, and \texttt{APOTEX INC} in
Canada).

To estimate how many medicines appear in multiple datasets based on exact matching, we build
a normalized key (\texttt{medicine\_id}) from \texttt{brandName} and \texttt{activeSubstances}
by converting text to upper case, removing accents and dosage strengths, standardizing
punctuation, and sorting multiple active substances. Let $F$, $E$ and $D$ denote the sets of
\texttt{medicine\_id} values in Drugs@FDA, EMA Medicines and the DPD. The pairwise exact
overlaps are $|F \cap E| = 428$, $|F \cap D| = 949$ and $|E \cap D| = 594$, yielding a
deduplicated union of $\left|(F \cap E) \cup (F \cap D) \cup (E \cap D)\right| = 1{,}407$
distinct medicines -- well above the required minimum of 1,000. (Note that this initial
exploratory estimate was calculated on raw tabular extracts prior to final EMA entity
consolidation; when evaluated via transitive graph clustering over the finalized XML datasets
in Section~\ref{subsec:clustering}, the exact baseline resolves to $1{,}372$ unified
multi-source medicine clusters).

However, exact matching misses many identical medicines because both identifying attributes
can vary across countries (Table~\ref{tab:difficult_match}):
\begin{itemize}
  \item \textbf{Active substance naming and salt forms:} The same substance can have different
  regional spellings (e.g., \texttt{CYCLOSPORINE} in the US and Canada vs.\
  \texttt{ciclosporin} in the EU for \textit{Verkazia}) or different chemical and salt
  representations (e.g., \texttt{GLYCOPYRROLATE} in the US, \texttt{glycopyrronium bromide}
  in the EU, and \texttt{GLYCOPYRRONIUM (GLYCOPYRRONIUM BROMIDE)} in Canada for
  \textit{Seebri}).
  \item \textbf{Brand name suffixes:} The same medicine may carry different device or
  formulation suffixes across markets, such as \texttt{SEEBRI NEOHALER} in Drugs@FDA versus
  \texttt{Seebri Breezhaler} in EMA Medicines and the DPD.
\end{itemize}

\begin{table}[htbp]
\centering
\caption{Examples of identical medicines with differing identifying attributes across sources}
\label{tab:difficult_match}
\footnotesize
\begin{tabularx}{\textwidth}{|l|l|X|}
\hline
\textbf{Dataset} & \textbf{Brand Name} & \textbf{Active Substance} \\
\hline
Drugs@FDA & \texttt{VERKAZIA} & \texttt{CYCLOSPORINE} \\
EMA Medicines & \texttt{Verkazia} & \texttt{ciclosporin} \\
Health Canada DPD & \texttt{VERKAZIA} & \texttt{CYCLOSPORINE} \\
\hline
Drugs@FDA & \texttt{SEEBRI NEOHALER} & \texttt{GLYCOPYRROLATE} \\
EMA Medicines & \texttt{Seebri Breezhaler} & \texttt{glycopyrronium bromide} \\
Health Canada DPD & \texttt{SEEBRI BREEZHALER} & \texttt{GLYCOPYRRONIUM (GLYCOPYRRONIUM BROMIDE)} \\
\hline
\end{tabularx}
\end{table}

% ============================================================
\section{Phase I: Data Selection and Translation}
% ============================================================

Phase~I translates three pharmaceutical datasets into a shared XML target
schema using Altova MapForce 2027. The outputs are validated XML files
consumed by Phases~II and~III.

% ------------------------------------------------------------
\subsection{Target Schema}
% ------------------------------------------------------------

The target schema (\texttt{medicines.xsd}) represents one medicine entity
per source record, identified by \texttt{brandName} and
\texttt{activeSubstances}, with 16 attributes and two metadata fields
(\texttt{id}, \texttt{source}). Three conventions apply across all
mappers: missing values are omitted rather than written as empty tags;
repeating values use typed wrapper elements (\texttt{activeSubstances},
\texttt{strengths}); and all dates follow ISO~8601 (YYYY-MM-DD).
Controlled vocabularies are enforced for \texttt{regulatoryStatus} (six
values: \textit{Authorised, Suspended, Pending, Withdrawn, Application
withdrawn, Refused}) and \texttt{productClass} (\textit{Human,
Veterinary, Disinfectant}).

% ------------------------------------------------------------
\subsection{Shared Helper Functions}
% ------------------------------------------------------------

A shared MapForce library (\texttt{wdi\_helpers.mfd}) provides ten
user-defined functions imported by all mappers
(Table~\ref{tab:udfs}). Lookup XML files cover 25 regulatory statuses,
227 dosage forms, 147 routes, and 202 strength units across all three
sources.

\begin{table}[htbp]
\centering
\caption{Shared UDFs used across mappers}
\label{tab:udfs}
\scriptsize
\begin{tabular}{|l|>{\raggedright\arraybackslash}p{7.5cm}|}
\hline
\textbf{Function} & \textbf{Purpose} \\
\hline
\texttt{makeMedicineId} & Constructs a stable entity ID from source prefix and record key \\
\texttt{getActiveSubstances} & Tokenizes substance string; supports INN fallback for blank primary field \\
\texttt{standardizeRegulatoryStatus} & Maps raw status values to the six controlled vocabulary entries \\
\texttt{standardizeProductClass} & Maps category values to Human / Veterinary / Disinfectant \\
\texttt{standardizeAtcCodes} & Combines human and vet ATC codes; deduplicates, sorts, corrects typos \\
\texttt{yesNoToBoolean} & Converts Yes/No strings to \texttt{xs:boolean}; omits element when false \\
\texttt{parseAuthorisationDate} & Converts DD/MM/YYYY to ISO~8601 YYYY-MM-DD \\
\texttt{stdDosageForm} / \texttt{dosageFormList} & Standardizes single or multi-value dosage form strings \\
\texttt{stdRoute} / \texttt{routeList} & Standardizes single or multi-value administration routes \\
\texttt{splitStrengths} & Tokenizes a DPD strength string into (value, unit) pairs \\
\hline
\end{tabular}
\end{table}

% ------------------------------------------------------------
\subsection{EMA Medicines Mapper}
% ------------------------------------------------------------

The EMA dataset originates from a flat table of 2,746 rows consolidated
into 2,648 unique medicine entities (\texttt{ema\_entities.csv}). Each
consolidated entity row maps to one output \texttt{<medicine>} element.
Because EMA does not publish strength or dosage form data in the register
export, \texttt{dosageForm}, \texttt{route}, and \texttt{strengths} are
structurally absent from EMA output. Table~\ref{tab:ema_mapping} lists all
field correspondences.

\begin{table}[htbp]
\centering
\caption{EMA Medicines field mapping}
\label{tab:ema_mapping}
\scriptsize
\begin{tabular}{|l|l|l|}
\hline
\textbf{Target field} & \textbf{Source column} & \textbf{Transformation} \\
\hline
\texttt{id}                   & medicine\_id                 & \texttt{makeMedicineId("EMA", id)} \\
\texttt{brandName}            & Name of medicine             & Direct copy \\
\texttt{activeSubstances}     & Active substance / INN       & \texttt{getActiveSubstances}; INN fallback \\
\texttt{activeSubstanceCount} & (derived)                    & \texttt{count(activeSubstances)} \\
\texttt{licenceHolder}        & MA holder                    & Direct copy \\
\texttt{regulatoryStatus}     & Medicine status              & \texttt{standardizeRegulatoryStatus} \\
\texttt{productClass}         & Category                     & \texttt{standardizeProductClass} \\
\texttt{atcCode}              & ATC code (human) + ATCvet    & \texttt{standardizeAtcCodes} \\
\texttt{generic} / \texttt{biosimilar} / \texttt{orphanMedicine} & Generic / Biosimilar / Orphan & \texttt{yesNoToBoolean} \\
\texttt{euAuthorisationDate}  & Marketing authorisation date & \texttt{parseAuthorisationDate} \\
\hline
\end{tabular}
\end{table}

The most intricate correspondence handles active substances. For 57 of
the 2,746 rows, both \texttt{Active substance} and the \texttt{INN /
common name} column are blank; because \texttt{activeSubstances} is
mandatory in the target schema, these rows are excluded via a row filter.
A further 28 rows have a blank primary field but a populated INN column:
the mapper applies an INN fallback, normalizing whitespace and selecting
the non-empty value. Commas and the word ``and'' are not treated as
separators, preserving multi-word biological descriptions. Consolidating
the 2,689 usable rows by normalized brand name and active substance set
yields \textbf{2,648 unique medicine entities} (\texttt{ema\_entities.csv}).

EMA maintains separate human and veterinary ATC code columns.
\texttt{standardizeAtcCodes} concatenates both, strips asterisks used
for provisional codes, drops the \textit{Not yet assigned} placeholder,
corrects a known EMA typo (\texttt{VO4D} $\rightarrow$ \texttt{V04D}),
then deduplicates, sorts, and joins codes with \texttt{"; "}. Date
conversion (\texttt{parseAuthorisationDate}) reorders DD/MM/YYYY to
YYYY-MM-DD. The three boolean flags are stored as Yes/No strings;
\texttt{yesNoToBoolean} converts them and omits the element when false,
consistent with the \texttt{minOccurs="0"} schema convention. The full
export passed XSD validation with \textbf{2,648 unique \texttt{EMA\_} IDs}
and zero duplicate or schema errors.

% ------------------------------------------------------------
\subsection{Drugs@FDA Mapper}
% ------------------------------------------------------------

The Drugs@FDA source consists of relational TSV tables in which multiple
product rows across applications (\texttt{ApplNo}) correspond to the same
medicine entity. Product records from \texttt{Products.txt} are joined to
\texttt{Applications.txt} on \texttt{ApplNo}, linked to marketing status
lookup values, and consolidated by normalized \texttt{brandName} (\texttt{DrugName})
and sorted \texttt{activeSubstances} (\texttt{ActiveIngredient}) into
\textbf{8,175 unique medicine entities}. Eleven orphan product rows lacking
a matching application and records with empty active ingredients (such as
\texttt{ApplNo=221337}, \texttt{MEDICAL AIR, USP}) are excluded during
preparation. Table~\ref{tab:fda_mapping} summarises the field mappings.

\begin{table}[htbp]
\centering
\caption{Drugs@FDA field mapping}
\label{tab:fda_mapping}
\scriptsize
\setlength{\tabcolsep}{4pt}
\begin{tabular}{|l|l|>{\raggedright\arraybackslash}p{5.5cm}|}
\hline
\textbf{Target field} & \textbf{Source column} & \textbf{Transformation} \\
\hline
\texttt{id}                   & medicine\_id                     & \texttt{makeMedicineId("FDA", id)} \\
\texttt{brandName}            & DrugName (plurality)             & Direct copy \\
\texttt{activeSubstances}     & ActiveIngredient                 & \texttt{getActiveSubstances}; separator \texttt{"; "} \\
\texttt{activeSubstanceCount} & (derived)                        & \texttt{count(activeSubstances)} \\
\texttt{licenceHolder}        & SponsorName (plurality)          & Direct copy \\
\texttt{regulatoryStatus}     & MarketingStatusID                & \texttt{standardizeRegulatoryStatus} \\
\texttt{productClass}         & (constant)                       & \texttt{"Human"} \\
\texttt{generic}              & ApplType (\texttt{ANDA}/\texttt{NDA}/\texttt{BLA}) & \texttt{true} if all \texttt{ANDA}; \texttt{false} if all \texttt{NDA}/\texttt{BLA}; omit if mixed \\
\texttt{dosageForm}           & Form (form component)            & \texttt{fdaFormPart} $\rightarrow$ \texttt{dosageFormList} \\
\texttt{route}                & Form (route component)           & \texttt{fdaRoutePart} $\rightarrow$ \texttt{routeList} \\
\texttt{strengths}            & Strength                         & Pre-split combination tokens $\rightarrow$ \texttt{splitStrengths} \\
\texttt{usApprovalDate}       & SubmissionStatusDate (\texttt{AP}) & Earliest non-empty approval date (ISO~8601) \\
\hline
\end{tabular}
\end{table}

Drugs@FDA stores dosage form and administration route jointly in
\texttt{Products.Form} (e.g.\ \texttt{"TABLET;ORAL"}); the helper functions
\texttt{fdaFormPart} and \texttt{fdaRoutePart} split this composite field
before lookup standardization. Combination strength strings such as
\texttt{600MG/90MG} are pre-split into individual mass tokens prior to
numeric value/unit extraction. Because Drugs@FDA covers only human drugs,
\texttt{productClass} maps to \textit{Human}. The \texttt{generic} flag is
derived from application types: \texttt{true} when all contributing
applications are \texttt{ANDA}, \texttt{false} when all are innovator
applications (\texttt{NDA} or \texttt{BLA}), and omitted for the
approximately 375 mixed entities. The resulting XML export
(\texttt{fda\_medicines.xml}) passed XSD and duplicate validation with
\textbf{8,175 unique \texttt{FDA\_} IDs}.

% ------------------------------------------------------------
\subsection{Health Canada DPD Mapper}
% ------------------------------------------------------------

The DPD source is a pre-consolidated CSV of \textbf{41,636 medicine
entities} derived from 58,137 product records across 12 Health Canada
tables (\texttt{dpd\_entities.csv}). Unlike the EMA flat table, each DPD
entity already aggregates multiple product records, so the mapper
operates on one row per entity without further consolidation. The DPD
source provides all 16 target attributes including \texttt{dosageForm},
\texttt{route}, and \texttt{strengths}, making it the most complete of
the three sources.

\begin{table}[htbp]
\centering
\caption{Health Canada DPD field mapping}
\label{tab:dpd_mapping}
\scriptsize
\begin{tabular}{|l|l|l|}
\hline
\textbf{Target field} & \textbf{Source column} & \textbf{Transformation} \\
\hline
\texttt{id}                   & medicine\_id          & \texttt{makeMedicineId("DPD", id)} \\
\texttt{brandName}            & brand\_name           & Direct copy \\
\texttt{activeSubstances}     & active\_substances    & \texttt{getActiveSubstances}; separator \texttt{"; "} \\
\texttt{activeSubstanceCount} & (derived)             & \texttt{count(activeSubstances)} \\
\texttt{licenceHolder}        & licence\_holder       & Direct copy \\
\texttt{regulatoryStatus}     & current\_status       & \texttt{standardizeRegulatoryStatus} \\
\texttt{productClass}         & product\_class        & \texttt{standardizeProductClass} \\
\texttt{atcCode}              & atc\_code             & \texttt{standardizeAtcCodes} \\
\texttt{biosimilar}           & biosimilar            & \texttt{yesNoToBoolean} \\
\texttt{dosageForm}           & dosage\_forms         & \texttt{dosageFormList} \\
\texttt{route}                & routes                & \texttt{routeList} \\
\texttt{strengths}            & strengths             & \texttt{splitStrengths} $\rightarrow$ (value, unit) pairs \\
\texttt{caStatusDate}         & status\_date          & Direct copy (ISO date) \\
\hline
\end{tabular}
\end{table}

The main mapping challenge in DPD is strength parsing. Each entity row
carries a semicolon-delimited strength string (e.g.\
\texttt{"10MG; 0.5\%"}); \texttt{splitStrengths} tokenizes this into
individual tokens, from which \texttt{strengthValue} extracts the numeric
component and \texttt{strengthUnit} normalizes the unit via
\texttt{lookups/strength\_unit.xml}. Tokens without a parseable numeric
value are filtered out before the \texttt{<strength>} wrapper is created,
and the wrapper itself is omitted when no valid tokens remain. Regulatory
status presents a similar aggregation challenge: DPD entities can carry
multiple historical statuses (e.g.\ \texttt{"MARKETED; CANCELLED POST
MARKET"}), so \texttt{standardizeRegulatoryStatus} selects the highest
active status according to the precedence hierarchy
$\text{Authorised} > \text{Suspended} > \cdots > \text{Refused}$.

The full export passed XSD validation with \textbf{41,636 records},
100,094 strength pairs across 41,574 medicines, and no schema errors.

% ------------------------------------------------------------
\subsection{Integrated XML Validation Summary}
% ------------------------------------------------------------

All three translated XML files in \texttt{data/4\_translated\_mapforce/}
were validated against \texttt{medicines.xsd} and checked for duplicate
IDs, duplicate attribute content, and duplicate normalized entity keys
using \texttt{check\_schema.py} (Table~\ref{tab:xml_validation}). Together,
the three sources provide \textbf{52,459 validated medicine entities} for
Phase~II identity resolution.

\begin{table}[htbp]
\centering
\caption{Validation summary of translated XML datasets}
\label{tab:xml_validation}
\scriptsize
\begin{tabular}{|l|r|r|r|}
\hline
\textbf{Validation check} & \textbf{\texttt{dpd\_medicines.xml}} & \textbf{\texttt{ema\_medicines.xml}} & \textbf{\texttt{fda\_medicines.xml}} \\
\hline
Records                       & 41,636 & 2,648 & 8,175 \\
Unique IDs                    & 41,636 & 2,648 & 8,175 \\
XSD Schema valid              & True   & True  & True  \\
Duplicate IDs                 & 0      & 0     & 0     \\
Duplicate content (excl.\ ID) & 0      & 0     & 0     \\
Duplicate normalized entities & 0      & 0     & 0     \\
\hline
\end{tabular}
\end{table}


% ============================================================
\section{Phase II: Identity Resolution}
% ============================================================

Phase~II resolves cross-jurisdictional medicine correspondences across
Health Canada DPD ($41{,}636$ entities), EMA Medicines ($2{,}648$ entities),
and Drugs@FDA ($8{,}175$ entities). Evaluating the full Cartesian product
across the three dataset pairs would require
$41{,}636 \times 2{,}648 + 41{,}636 \times 8{,}175 + 2{,}648 \times 8{,}175 = 472{,}273{,}828$
pairwise comparisons, motivating blocking, supervised and zero-shot matchers,
and transitive connected-components clustering.

% ------------------------------------------------------------
\subsection{Gold-Standard Dataset Construction and Stratified Split}
% ------------------------------------------------------------

To construct a representative benchmark that tests both standard matches
and adversarial pharmaceutical corner cases, candidate pairs for manual
annotation were sampled in two balanced pools using
\texttt{gold\_standard\_pair\_gen.py}:
\begin{itemize}
  \item \textbf{Ordinary Pool (\texttt{hard=0}):} Contains candidate pairs
  with exact or high-similarity normalized brand and active substance keys
  alongside random negative pairs.
  \item \textbf{Hard Pool (\texttt{hard=1}):} Targets domain-specific
  corner cases where string similarity is misleading: (i)~identical brand
  names with different active substances or combination vs.\ single-agent
  formulas, (ii)~identical active substances marketed by competing generic
  manufacturers, and (iii)~same medicine entities with regional INN salt
  variations or device/formulation suffixes.
\end{itemize}

Three team annotators manually reviewed all candidate batches across
\texttt{DPD/EMA} ($1{,}000$ pairs), \texttt{DPD/FDA} ($769$ pairs), and
\texttt{EMA/FDA} ($1{,}729$ pairs), yielding \textbf{3,498 unique reviewed
pairs} with zero duplicate rows, zero conflicting labels, and zero unlabeled
records. Two records were labeled as a match ($\texttt{match}=1$) if and only
if they share the same pharmacological active substance set (accounting for
salt forms, hydration states, and regional INN spellings) and represent the
same commercial medicine entity (including regional subsidiaries of the same
corporate group, while excluding unrelated competing generic manufacturers).

Using \texttt{gold\_standard\_combine\_split.py}, the 3,498 unified pairs
were partitioned into an 85\% development split (\texttt{gs\_pair\_train\_val.csv},
\textbf{2,973 pairs}) and a 15\% held-out test split (\texttt{gs\_pair\_test.csv},
\textbf{525 pairs}) stratified jointly by $(\texttt{source\_pair}, \texttt{hard}, \texttt{match})$
with fixed seed 42 (Table~\ref{tab:gs_split}).

\begin{table}[htbp]
\centering
\caption{Stratified gold-standard dataset distribution across source pairs, difficulty pools, and partitions}
\label{tab:gs_split}
\scriptsize
\setlength{\tabcolsep}{4pt}
\begin{tabular}{|l|l|l|r|r|r|}
\hline
\textbf{Source Pair} & \textbf{Difficulty Pool} & \textbf{Label} & \textbf{Whole} & \textbf{Train/Val (85\%)} & \textbf{Test (15\%)} \\
\hline
DPD / EMA & Ordinary (\texttt{hard=0}) & Match (1)     &   198 &   168 &  30 \\
DPD / EMA & Ordinary (\texttt{hard=0}) & Non-match (0) &   302 &   257 &  45 \\
DPD / EMA & Hard (\texttt{hard=1})     & Match (1)     &   217 &   184 &  33 \\
DPD / EMA & Hard (\texttt{hard=1})     & Non-match (0) &   283 &   241 &  42 \\
\hline
DPD / FDA & Ordinary (\texttt{hard=0}) & Match (1)     &   122 &   104 &  18 \\
DPD / FDA & Ordinary (\texttt{hard=0}) & Non-match (0) &   364 &   309 &  55 \\
DPD / FDA & Hard (\texttt{hard=1})     & Match (1)     &   126 &   107 &  19 \\
DPD / FDA & Hard (\texttt{hard=1})     & Non-match (0) &   157 &   133 &  24 \\
\hline
EMA / FDA & Ordinary (\texttt{hard=0}) & Match (1)     &   282 &   240 &  42 \\
EMA / FDA & Ordinary (\texttt{hard=0}) & Non-match (0) &   670 &   570 & 100 \\
EMA / FDA & Hard (\texttt{hard=1})     & Match (1)     &   292 &   248 &  44 \\
EMA / FDA & Hard (\texttt{hard=1})     & Non-match (0) &   485 &   412 &  73 \\
\hline
\textbf{Total} & \textbf{All Pools} & \textbf{All Labels} & \textbf{3,498} & \textbf{2,973} & \textbf{525} \\
\hline
\end{tabular}
\end{table}

% ------------------------------------------------------------
\subsection{Blocking and Candidate Generation}
% ------------------------------------------------------------

The unreduced search space across the three translated datasets encompasses
$41{,}636 \times 2{,}648 = 110{,}252{,}128$ pairs for DPD--EMA,
$41{,}636 \times 8{,}175 = 340{,}374{,}300$ pairs for DPD--FDA, and
$2{,}648 \times 8{,}175 = 21{,}647{,}400$ pairs for EMA--FDA, yielding a total
Cartesian product of \textbf{472,273,828 possible comparisons}. Performing
detailed attribute matching or transformer inference across nearly half a billion
pairs is computationally prohibitive. The goal of blocking is to filter out
obvious non-matches while maintaining high completeness on true correspondences.
Performance is governed by two complementary metrics:
\begin{itemize}
  \item \textbf{Reduction Ratio (RR):} $\text{RR} = 1 - \frac{|\mathcal{C}|}{|A \times B|}$,
  measuring the relative reduction in pairwise comparisons (target: $>99.5\%$).
  \item \textbf{Pairs Completeness (PC):} $\text{PC} = \frac{|\mathcal{C} \cap \mathcal{M}|}{|\mathcal{M}|}$,
  the fraction of true matches $\mathcal{M}$ retained in candidate shortlist $\mathcal{C}$
  (target: $>97.0\%$ per source pair).
\end{itemize}

\paragraph{Lexical Blocking Baselines.}
We initially benchmarked standard lexical blocking rules implemented via PyDI
(\texttt{StandardBlocker} and \texttt{TokenBlocker}) on the development split
(\texttt{gs\_pair\_train\_val.csv}, $1{,}038$ unique eligible positive matches
after EMA ID alignment and deduplication; Table~\ref{tab:blocking_development}):
\begin{enumerate}
  \item \textbf{Brand Prefix-4:} Compares the first four alphanumeric characters
  of normalized brand names. While achieving high reduction ratios ($>99.96\%$),
  it missed 39 of the 352 DPD--EMA matches ($\text{PC} = 88.92\%$) due to regional
  brand prefixes (e.g., \texttt{TEVA-}, \texttt{APO-}, or formulation prefixes).
  \item \textbf{Brand Trigrams (capped at 1,000):} Tokenizes brand names into
  overlapping 3-character shingles. To prevent Cartesian explosions on frequent
  pharmaceutical fragments (e.g., \texttt{"TAB"}, \texttt{"SOL"}), candidate
  generation was capped at $a \times b \le 1{,}000$ co-occurrences per shingle.
  However, trigrams alone missed 39 matches on DPD--FDA ($\text{PC} = 81.52\%$).
  \item \textbf{Prefix-4 $\cup$ Trigrams:} Combining both lexical rules improved
  completeness to $\ge 98.30\%$ on all pairs, but generated $1{,}369{,}537$
  candidates ($663{,}705$ on DPD--FDA alone), representing a substantial
  downstream processing burden.
\end{enumerate}

\paragraph{Dense Semantic Embeddings (Qwen3-Embedding-0.6B).}
To capture semantic equivalence across international nonproprietary names (INNs),
salt variations, and brand synonyms, we evaluated dense vector blocking using
\texttt{Qwen/Qwen3-Embedding-0.6B} ($1{,}024$-dimensional representations).
For each medicine entity, the normalized brand name and alphabetically sorted
active substances were concatenated and encoded without fine-tuning.
Cosine nearest-neighbor search was executed from the smaller source into the
larger source. A shortlist of the top-5 nearest neighbors per medicine produced
only \textbf{67,355 total candidates} while retaining 1,021 of 1,038 development
positives (\textbf{98.36\% pooled completeness}; Table~\ref{tab:blocking_development}).

\begin{table}[htbp]
\centering
\caption{Comparison of lexical and dense embedding blocking methods on the development set ($1{,}038$ positive matches)}
\label{tab:blocking_development}
\scriptsize
\setlength{\tabcolsep}{2.5pt}
\begin{tabular}{|l|l|r|r|r|r|r|}
\hline
\textbf{Source Pair} & \textbf{Blocking Method} & \textbf{Candidates} & \textbf{Matches} & \textbf{Retained} & \textbf{Completeness} & \textbf{Red.\ Ratio} \\
\hline
DPD--EMA & Brand Prefix-4             &  16,398 & 352 & 313 & 88.92\% & 99.9851\% \\
DPD--EMA & Brand Trigrams (cap 1,000) & 431,353 & 352 & 323 & 91.76\% & 99.6088\% \\
DPD--EMA & Prefix-4 $\cup$ Trigrams   & 443,207 & 352 & 346 & 98.30\% & 99.5980\% \\
DPD--EMA & Qwen Top-5                 &  13,240 & 352 & 346 & 98.30\% & 99.9880\% \\
DPD--EMA & \textbf{Qwen Top-5 $\cup$ Prefix-4} & \textbf{27,966} & \textbf{352} & \textbf{350} & \textbf{99.43\%} & \textbf{99.9746\%} \\
\hline
DPD--FDA & Brand Prefix-4             & 133,069 & 211 & 210 & 99.53\% & 99.9609\% \\
DPD--FDA & Brand Trigrams (cap 1,000) & 544,036 & 211 & 172 & 81.52\% & 99.8402\% \\
DPD--FDA & Prefix-4 $\cup$ Trigrams   & 663,705 & 211 & 210 & 99.53\% & 99.8050\% \\
DPD--FDA & Qwen Top-5                 &  40,875 & 211 & 210 & 99.53\% & 99.9880\% \\
DPD--FDA & \textbf{Qwen Top-5 $\cup$ Prefix-4} & \textbf{162,930} & \textbf{211} & \textbf{211} & \textbf{100.00\%} & \textbf{99.9521\%} \\
\hline
EMA--FDA & Brand Prefix-4             &   4,456 & 475 & 468 & 98.53\% & 99.9794\% \\
EMA--FDA & Brand Trigrams (cap 1,000) & 261,601 & 475 & 472 & 99.37\% & 98.7915\% \\
EMA--FDA & Prefix-4 $\cup$ Trigrams   & 262,625 & 475 & 474 & 99.79\% & 98.7868\% \\
EMA--FDA & Qwen Top-5                 &  13,240 & 475 & 465 & 97.89\% & 99.9388\% \\
EMA--FDA & \textbf{Qwen Top-5 $\cup$ Prefix-4} & \textbf{15,934} & \textbf{475} & \textbf{474} & \textbf{99.79\%} & \textbf{99.9264\%} \\
\hline
\end{tabular}
\end{table}

\paragraph{Frozen Evaluation and Final Matcher Handoff.}
When Qwen Top-5 was evaluated on the held-out test set ($185$ known matches),
it achieved $98.41\%$ on DPD--EMA and $97.30\%$ on DPD--FDA, but dropped to
$92.94\%$ ($79/85$) on EMA--FDA ($95.68\%$ pooled completeness), failing the
per-source $97\%$ target. Diagnostic error inspection revealed that missed
EMA--FDA pairs shared brand prefixes that lexical prefix matching reliably captures.

To guarantee compliance with the project's $\ge 97\%$ completeness requirement
across all three pairs, the final handoff rule was defined as the union:
$$\mathcal{C}_{\text{final}} = \mathcal{C}_{\text{Qwen top-5}} \cup \mathcal{C}_{\text{Prefix-4}}$$
As shown in Table~\ref{tab:blocking_candidates}, this hybrid rule recovers 6
additional test matches, bringing EMA--FDA to \textbf{97.65\% completeness}
and achieving \textbf{100.00\% completeness} on DPD--EMA and DPD--FDA
(\textbf{98.92\% pooled completeness}, $183/185$). The candidate shortlist
totals \textbf{206,830 unique pairs}, compressing the search space by
\textbf{99.9562\%} (over $2{,}280\times$ reduction). All candidates were
persisted to \texttt{data/7\_blocking\_candidates/} with cryptographic SHA-256
checksums verified in \texttt{manifest.json}.

\begin{table}[htbp]
\centering
\caption{Final candidate handoff summary and held-out test set evaluation ($185$ known matches)}
\label{tab:blocking_candidates}
\scriptsize
\setlength{\tabcolsep}{3pt}
\begin{tabular}{|l|r|r|r|r|r|r|}
\hline
\textbf{Source Pair} & \textbf{Cartesian Space} & \textbf{Candidates} & \textbf{Red.\ Ratio} & \textbf{Test Matches} & \textbf{Retained} & \textbf{Completeness} \\
\hline
DPD--EMA & 110,252,128 &  27,966 & 99.9746\% & 63 & 63 & \textbf{100.00\%} \\
DPD--FDA & 340,374,300 & 162,930 & 99.9521\% & 37 & 37 & \textbf{100.00\%} \\
EMA--FDA &  21,647,400 &  15,934 & 99.9264\% & 85 & 83 &  \textbf{97.65\%} \\
\hline
\textbf{Pooled} & \textbf{472,273,828} & \textbf{206,830} & \textbf{99.9562\%} & \textbf{185} & \textbf{183} & \textbf{98.92\%} \\
\hline
\end{tabular}
\end{table}

% ------------------------------------------------------------
\subsection{Rule-Based, Machine Learning, and Fine-Tuned PLM Matchers}
% ------------------------------------------------------------

% TODO (Hai, Trang, Özlem): Insert comparator definitions, Rule-Based
% matcher tuning, supervised ML classifiers (Decision Tree, Random Forest,
% SVM, Logistic Regression), and fine-tuned PLM evaluation results.

% ------------------------------------------------------------
\subsection{Zero-Shot LLM Matcher Benchmark}
% ------------------------------------------------------------

To evaluate whether a pretrained Large Language Model can resolve
cross-jurisdictional pharmaceutical entities without task-specific training
examples, we implemented a strictly blind Zero-Shot LLM matcher
(\texttt{src/identity\_resolution/llm\_zero\_shot\_matcher.py}) using
Google's \texttt{gemini-3.8-flash} model via the official \texttt{google-genai}
SDK ($\texttt{temperature}=0.0$, structured JSON schema output enforcing
integer $\texttt{match} \in \{0, 1\}$ and clinical rationale string
\texttt{reason}).

\paragraph{Decontaminated Zero-Shot Prompt Architecture.}
Each test pair in \texttt{gs\_pair\_test.csv} ($n = 525$) is formatted
with four record attributes per jurisdiction: \texttt{brandName},
\texttt{activeSubstances}, \texttt{licenceHolder}, and \texttt{dosageForm}
(with missing values formatted as \texttt{"Not specified"}). The system
instruction encodes three general domain rules without exposing any
test-set drug names or few-shot examples:
\begin{enumerate}
  \item \textbf{Active Substances:} Ingredients must represent the same
  pharmacological moiety; salt forms (e.g.\ hydrochloride, sodium, maleate,
  besylate), hydration states, and INN vs.\ USAN spellings (e.g.\
  ciclosporin vs.\ cyclosporine) are treated as equivalent, whereas
  single-substance vs.\ multi-substance combination products are non-matches.
  \item \textbf{Brand Name and Licence Holder:} Identical brands or minor
  formulation suffixes from regional subsidiaries of the same corporate
  group are matches, whereas generic products from unrelated competing
  manufacturers are non-matches.
  \item \textbf{Dosage Form Variations:} Minor dosage form differences do
  not invalidate a match when brand name, active substance set, and
  corporate lineage agree.
\end{enumerate}

\paragraph{Comparison Against the Naive Exact-Match Baseline.}
We benchmark the Zero-Shot LLM matcher against a deterministic
\textbf{Naive Exact-Match Baseline} that predicts $\texttt{match}=1$ if
and only if normalized \texttt{brandName1} equals \texttt{brandName2} and
their normalized sets of semicolon-delimited \texttt{activeSubstances}
are identical. Table~\ref{tab:llm_vs_baseline} reports overall test-set
performance across all 525 held-out pairs.

\begin{table}[htbp]
\centering
\caption{Held-out test set evaluation ($n = 525$ pairs): Naive Exact-Match Baseline vs.\ Zero-Shot LLM (\texttt{gemini-3.8-flash})}
\label{tab:llm_vs_baseline}
\scriptsize
\setlength{\tabcolsep}{4pt}
\begin{tabular}{|l|r|r|r|}
\hline
\textbf{Metric} & \textbf{Exact-Match Baseline} & \textbf{Zero-Shot LLM (\texttt{gemini-3.8-flash})} & \textbf{Improvement ($\Delta$)} \\
\hline
\textbf{Macro F1 Score}       & 0.7307 (73.07\%) & \textbf{0.9485 (94.85\%)} & \textbf{+21.78 pp} \\
\textbf{Match Class F1 (1)}   & 0.6122 (61.22\%) & \textbf{0.9344 (93.44\%)} & \textbf{+32.21 pp} \\
\textbf{Non-Match Class F1 (0)} & 0.8492 (84.92\%) & \textbf{0.9626 (96.26\%)} & \textbf{+11.34 pp} \\
\textbf{Overall Accuracy}     & 0.7829 (78.29\%) & \textbf{0.9524 (95.24\%)} & \textbf{+16.95 pp} \\
Macro Precision               & 0.8016 (80.16\%) & 0.9443 (94.43\%)          & +14.27 pp \\
Macro Recall                  & 0.7154 (71.54\%) & 0.9534 (95.34\%)          & +23.80 pp \\
Match Precision (Class 1)     & 0.8333 (83.33\%) & 0.9128 (91.28\%)          & +7.95 pp \\
Match Recall (Class 1)        & 0.4839 (48.39\%) & 0.9570 (95.70\%)          & +47.31 pp \\
\hline
True Positives (TP)           & 90               & \textbf{178}              & +88 matches found \\
False Negatives (FN)          & 96               & \textbf{8}                & $-88$ missed matches \\
False Positives (FP)          & 18               & \textbf{17}               & $-1$ false alarm \\
True Negatives (TN)           & 321              & \textbf{322}              & +1 \\
\hline
\end{tabular}
\end{table}

The Naive Exact-Match Baseline achieves a Macro F1 of $73.07\%$ and misses
more than half of all true matches ($\text{Recall}_1 = 48.39\%$, $\text{FN} = 96$)
because regional salt modifiers and brand suffixes break exact string
equality. By contrast, \texttt{gemini-3.8-flash} achieves \textbf{94.85\%
Macro F1} ($95.24\%$ accuracy), recovering 88 additional true matches and
reducing false negatives by $91.7\%$ (from 96 to 8).

\paragraph{Performance by Difficulty Pool and Source Jurisdiction Pair.}
Table~\ref{tab:llm_breakdowns} disaggregates test-set performance across
difficulty pools and regulatory pair combinations. On the adversarial
\textbf{Hard Pool} ($\texttt{hard}=1$, $235$ pairs), the Exact-Match
Baseline fails completely on positive matches ($\text{Match F1} = 0.0000$,
$\text{Macro F1} = 0.3717$), whereas the Zero-Shot LLM maintains
\textbf{0.9296 Macro F1} ($0.9167$ Match F1, $93.19\%$ accuracy). Across
source pairs, the LLM achieves $0.9864$ Macro F1 on \texttt{DPD/EMA},
$0.9704$ on \texttt{DPD/FDA}, and $0.9144$ on \texttt{EMA/FDA}.

\begin{table}[htbp]
\centering
\caption{Zero-Shot LLM vs.\ Naive Baseline breakdown by difficulty pool and source jurisdiction pair}
\label{tab:llm_breakdowns}
\scriptsize
\setlength{\tabcolsep}{4pt}
\begin{tabular}{|l|r|r|r|r|r|r|}
\hline
\textbf{Subgroup} & \textbf{Pairs} & \textbf{Base.\ Macro F1} & \textbf{Base.\ Match F1} & \textbf{LLM Acc.} & \textbf{LLM Macro F1} & \textbf{LLM Match F1} \\
\hline
\multicolumn{7}{|l|}{\textit{Breakdown by Difficulty Pool}} \\
\hline
Ordinary Pool (\texttt{hard=0}) & 290 & 0.9310 & 0.9091 & 96.90\% & \textbf{0.9647} & 0.9524 \\
Hard Pool (\texttt{hard=1})     & 235 & 0.3717 & 0.0000 & 93.19\% & \textbf{0.9296} & 0.9167 \\
\hline
\multicolumn{7}{|l|}{\textit{Breakdown by Source Jurisdiction Pair}} \\
\hline
DPD / EMA (Canada $\leftrightarrow$ EU) & 150 & 0.7366 & 0.6383 & 98.67\% & \textbf{0.9864} & 0.9844 \\
DPD / FDA (Canada $\leftrightarrow$ US) & 116 & 0.6979 & 0.5625 & 97.41\% & \textbf{0.9704} & 0.9600 \\
EMA / FDA (EU $\leftrightarrow$ US)     & 259 & 0.7408 & 0.6176 & 92.28\% & \textbf{0.9144} & 0.8876 \\
\hline
\end{tabular}
\end{table}

\paragraph{Qualitative Clinical Error Analysis.}
Inspecting all 25 test-set discrepancies ($17$ False Positives and $8$
False Negatives) reveals four domain patterns:
\begin{enumerate}
  \item \textbf{Cross-Border Licensing and Corporate Acquisitions (14 FP):}
  In 14 pairs with identical brand names and active substances (e.g.\
  \textit{Xolair} marketed by Novartis in the EU and Genentech in the US;
  \textit{Doptelet} by SOBI vs.\ wholly owned subsidiary AkaRx; \textit{Tekturna}
  by Novartis vs.\ acquirer LXO; \textit{Nevanac} by Novartis vs.\ Harrow Eye;
  \textit{Adasuve} by Ferrer vs.\ Nova Pneuma; \textit{Zurampic} by
  Gr\"{u}nenthal vs.\ Ironwood), the LLM leveraged pretrained domain
  knowledge of licensing partnerships to predict $\texttt{match}=1$, whereas
  annotators labeled them $0$ due to differing surface manufacturer names
  (plus \textit{Tevimbra} by BeOne Medicines Germany vs.\ Ireland, which is
  a gold-standard annotation error).
  \item \textbf{Formulation Suffix Variants (3 FP, 1 FN):}
  Modified delivery formulations from the same innovator (\textit{Natpar}
  vs.\ \textit{Natpara}, \textit{Alphagan} vs.\ \textit{Alphagan P}, and
  \textit{Afinitor Disperz} vs.\ \textit{Afinitor} predicted $1$ vs.\ GS $0$;
  \textit{Tobi Podhaler} dry powder inhaler vs.\ \textit{TOBI} nebulizer
  solution predicted $0$ vs.\ GS $1$).
  \item \textbf{Regional Co-Brands, Prodrugs, and Co-Formulations (3 FN):}
  \textit{Copalia HCT} vs.\ \textit{Exforge HCT} (distinct EU co-marketing
  brand names from Novartis), \textit{Ivemend} (\texttt{fosaprepitant}
  intravenous prodrug) vs.\ \textit{Emend} (\texttt{aprepitant} oral parent
  drug), and \textit{Herceptin} (\texttt{trastuzumab}) vs.\ \textit{Herceptin
  Hylecta} (\texttt{trastuzumab and hyaluronidase}).
  \item \textbf{Gold-Standard Label Noise Corrected by the LLM (4 FN):}
  Four pairs labeled $1$ in the manual gold standard are genuine
  non-matches where the LLM's $0$ prediction is clinically correct:
  \texttt{AURO-LAMIVUDINE/ZIDOVUDINE} (dual therapy) vs.\
  \texttt{EFAVIRENZ; LAMIVUDINE; ZIDOVUDINE} (triple therapy), and three
  competing generic pairs (\textit{Repaglinide Accord} vs.\ Aurobindo,
  \textit{Pemetrexed Krka} vs.\ Shilpa, and \textit{Tadalafil Lilly} vs.\
  Rising).
\end{enumerate}

% ------------------------------------------------------------
\subsection{Tri-Dataset Connected-Components Clustering}
\label{subsec:clustering}
% ------------------------------------------------------------

Pairwise entity resolution produces binary correspondence sets across
$\texttt{DPD}\leftrightarrow\texttt{EMA}$, $\texttt{DPD}\leftrightarrow\texttt{FDA}$,
and $\texttt{EMA}\leftrightarrow\texttt{FDA}$. To consolidate these pairwise
links into unified multi-source medicine entities for Phase~III fusion, we
implemented graph-based transitive closure in
\texttt{src/identity\_resolution/clustering.py} using NetworkX connected
components. Let $G = (V, E)$ be an undirected graph where vertices $V$ are
source-prefixed entity IDs ($\texttt{DPD\_*}$, $\texttt{EMA\_*}$, $\texttt{FDA\_*}$)
and edges $E$ are positive pairwise match correspondences. Each connected
component $C_k \subseteq V$ forms a unified medicine cluster assigned a
deterministic identifier (\texttt{CLUSTER\_00001}, \dots).

Executing \texttt{clustering.py} on the exact normalized baseline
correspondences across the three validated XML datasets produces
\textbf{1,372 cross-dataset medicine clusters} spanning at least two
jurisdictions, including \textbf{283 three-way clusters} present in
Canada, the EU, and the US simultaneously (Table~\ref{tab:clustering_profile}).
This confirms that the integrated dataset exceeds the course threshold of
$1{,}000$ overlapping entities even before incorporating fuzzy and learned
matcher correspondences. This count refines the preliminary exploratory
profiling estimate of $1{,}407$ reported in Section~\ref{subsec:entity_id}:
whereas the exploratory estimate counted distinct synthetic keys across
raw pre-consolidation CSV extracts, connected-components clustering
operates directly on the deduplicated XML target datasets (where EMA was
consolidated into $2{,}648$ unique entities) and applies transitive graph
closure to unify pairwise correspondences into coherent multi-source entity clusters.

Furthermore, when clustering is executed over the full candidate pool
generated by the hybrid blocker ($\mathcal{C}_{\text{Qwen top-5}} \cup \mathcal{C}_{\text{Prefix-4}}$,
$206{,}830$ candidate pairs) using our hierarchical matcher with salt-form
normalization and deterministic greedy $1{:}1$ conflict resolution
(\texttt{--enforce-one-to-one}), the yield of cross-dataset medicines
increases to \textbf{2,343 clusters} ($+70.8\%$ gain over the exact baseline),
with three-way international clusters increasing to \textbf{502} ($+77.4\%$ gain).

\begin{table}[htbp]
\centering
\caption{Connected-components clustering profile: Group size distribution and jurisdictional overlap breakdown (Baseline Exact vs.\ Candidate Matches with Greedy 1:1)}
\label{tab:clustering_profile}
\scriptsize
\setlength{\tabcolsep}{4pt}
\begin{tabular}{|l|r|r|r|r|}
\hline
\textbf{Metric / Category} & \textbf{Exact Baseline} & \textbf{Share} & \textbf{Candidate Matches} & \textbf{Share} \\
\hline
\multicolumn{5}{|l|}{\textit{\textbf{Group Size Distribution ($k$)}}} \\
\hline
Pairwise ($k = 2$)            & 1,089 &  79.37\% & 1,841 &  78.57\% \\
Tri-Jurisdictional ($k = 3$)  &   283 &  20.63\% &   497 &  21.21\% \\
Size $k = 4$                  &     0 &   0.00\% &     4 &   0.17\% \\
Size $k = 5$                  &     0 &   0.00\% &     1 &   0.04\% \\
\hline
\textbf{Total Multi-Source Clusters} & \textbf{1,372} & \textbf{100.00\%} & \textbf{2,343} & \textbf{100.00\%} \\
Total Entities Clustered      & 3,027 & ---      & 5,194 & ---      \\
Intra-Source Variations ($k > 3$) & 0 & 0.00\%   &     5 &   0.21\% \\
\hline
\multicolumn{5}{|l|}{\textit{\textbf{Jurisdictional Overlap Breakdown ($\ge 2$ Sources)}}} \\
\hline
All 3 Sources (\texttt{DPD + EMA + FDA}) &   283 &  20.63\% &   502 &  21.43\% \\
Canada \& US Only (\texttt{DPD + FDA})   &   629 &  45.85\% & 1,388 &  59.24\% \\
Canada \& EU Only (\texttt{DPD + EMA})   &   311 &  22.67\% &   314 &  13.40\% \\
EU \& US Only (\texttt{EMA + FDA})       &   149 &  10.86\% &   139 &   5.93\% \\
\hline
\textbf{Total Overlapping Entities}      & \textbf{1,372} & \textbf{100.00\%} & \textbf{2,343} & \textbf{100.00\%} \\
\hline
\end{tabular}
\end{table}

\paragraph{Cluster / Group Size Distribution.}
A crucial quality metric for identity resolution is the group size distribution resulting from transitive closure over match correspondences. Table~\ref{tab:clustering_profile} summarizes the resulting cluster sizes for both the Exact Baseline and the Candidate Correspondences.
\begin{itemize}
  \item \textbf{Exact Baseline:} Transitive closure unifies $3{,}027$ medicine records into $1{,}372$ clusters (mean size $2.21$). Exactly $1{,}089$ clusters ($79.37\%$) are pairwise ($k=2$, linking two jurisdictions), and $283$ clusters ($20.63\%$) are triplets ($k=3$, linking Canada, the EU, and the US). No clusters exceed size 3 ($100.0\%$ pure $1{:}1$ and $1{:}1{:}1$ equivalence classes with zero intra-source conflicts).
  \item \textbf{Candidate Matches (Greedy 1:1):} Transitive closure consolidates $5{,}194$ entity records into $2{,}343$ multi-source clusters (mean size $2.22$). Pairwise clusters account for $1{,}841$ clusters ($78.57\%$), and tri-jurisdictional triplets account for $497$ clusters ($21.21\%$). Only 5 clusters ($0.21\%$) contain more than three records (four of size 4, one of size 5), demonstrating that $99.79\%$ of resolved candidate clusters are strictly duplicate-free single-entry-per-source components.
\end{itemize}

\paragraph{Analysis of Clusters with Size $> 3$.}
Following course guidelines, clusters containing multiple entities from the same source jurisdiction were manually audited to assess whether they originate from false-positive transitive chaining or legitimate intra-source duplicates:
\begin{enumerate}
  \item \texttt{CLUSTER\_00001} ($k=5$): Fixed-dose combination \textit{Irbesartan and Hydrochlorothiazide}. Both EMA and FDA contain separate marketing authorizations for competing generic manufacturers (EMA has BMS and Teva licences; FDA has two distinct ANDA applications). Transitive closure links these duplicate generic authorisations to the single Canadian DPD entry, yielding 5 records.
  \item \texttt{CLUSTER\_00002} ($k=4$): \textit{Temozolomide}, an oncology therapeutic. Health Canada DPD contains two separate active DIN approvals from competing generic sponsors (\textit{ACH Temozolomide} and \textit{Temozolomide Capsules}) linking to single EMA and FDA records.
  \item \texttt{CLUSTER\_00003} ($k=4$): \textit{Sitagliptin + Metformin Hydrochloride}. DPD contains two distinct generic DIN approvals (\textit{Apo-Sitagliptin} and \textit{Reddy-Sitagliptin}) matching the corresponding EMA and FDA entries.
  \item \texttt{CLUSTER\_00004} ($k=4$): \textit{Emtricitabine + Tenofovir Disoproxil}. DPD registers two distinct generic approvals (\textit{Jamp} and \textit{Mylan}) clustering with the international reference entries.
  \item \texttt{CLUSTER\_00005} ($k=4$): \textit{Koselugo} (\textit{selumetinib}). US Drugs@FDA lists two separate application records reflecting base substance (\textit{selumetinib}) versus sulfate salt (\textit{selumetinib sulfate}) representations, both clustering with the Canadian and European entries.
\end{enumerate}
In all five cases, the expanded group size reflects legitimate upstream intra-source generic duplicates or multi-application salt variations in the regulatory registers rather than false-positive cross-substance matching errors.

% ============================================================
\section{Phase III: Data Fusion}
\label{sec:data_fusion}
% ============================================================

Phase~III integrates attributes across clustered entities, resolves conflicting values, and measures the quality profile of the consolidated pharmaceutical dataset. In accordance with the course report guidelines, this section documents the experimental results across six dimensions:

\subsection{Provenance Tracking}
Each fused entity and individual attribute value is enriched with provenance metadata, recording the source regulatory registry (\texttt{Drugs@FDA}, \texttt{EMA}, \texttt{DPD}), original national identifiers, and dataset release dates (Drugs@FDA 2026, EMA 2026, DPD 2026) to provide full lineage and auditability.

\subsection{Dataset Density and Consistency Assessment}
We evaluate the density of each input dataset prior to fusion alongside the merged dataset across the 16 target attributes. Intra-source and cross-source attribute consistency is analyzed to identify systematic discrepancies across jurisdictions (e.g., brand naming suffixes, salt representations, and regulatory status categories).

\subsection{Fusion Evaluation Benchmark Set}
To rigorously assess conflict resolution quality, an evaluation benchmark set containing $\ge 15$ multi-source medicine entities with $\ge 4$ attributes per entity is established. Ground-truth values are independently validated against external regulatory monographs and clinical pharmacopeia (e.g., DailyMed, WHO INN register), ensuring unambiguous single-record references without evaluation ambiguities.

\subsection{Conflict Resolution Functions and Custom Rules}
We systematically experiment with standard conflict resolution strategies (Voting, Longest String, Most Trusted Source based on jurisdiction data quality scores, and Union for multi-valued active substances and strengths) as well as custom pharmaceutical domain heuristics (e.g., INN-preference for substance nomenclature and hierarchical dosage form mapping).

\subsection{Conflict Resolution Accuracy Comparison}
The performance of different resolution functions is comparatively evaluated against the benchmark set. A comparative table reports attribute-level accuracy across alternative functions, identifying the optimal strategy for each target attribute.

\subsection{Remaining Error Analysis}
An error audit is performed on the remaining discrepancies following the application of the best resolution functions, categorizing residual errors into regulatory classification differences, temporal update lags, and typographical reporting variants.

% ============================================================
\section{Conclusion}
\label{sec:conclusion}
% ============================================================

In accordance with the project report requirements, the overall integration results are summarized by addressing three core evaluation questions:
\begin{enumerate}
  \item \textbf{Additional Entities Added:} The total number of unique medicine entities gained beyond the largest individual input dataset (Health Canada DPD with $41{,}636$ entities), quantifying the expanded international coverage achieved through tri-dataset alignment.
  \item \textbf{Data Density Increase:} The net increase in attribute density achieved in the fused dataset relative to the highest-density input source, demonstrating the complementary informational yield of multi-jurisdiction merging.
  \item \textbf{Overall Dataset Accuracy:} The end-to-end accuracy of the final fused dataset as measured against the independent ground-truth evaluation benchmark.
\end{enumerate}




% ============================================================
% SOURCES
% ============================================================

\section*{Sources}
\addcontentsline{toc}{section}{Sources}

\begin{itemize}
    \item U.S. Food and Drug Administration,
    \emph{Drugs@FDA Data Files}.\\
    \url{https://www.fda.gov/drugs/drug-approvals-and-databases/drugsfda-data-files}

    \item European Medicines Agency,
    \emph{Download medicine data}.\\
    \url{https://www.ema.europa.eu/en/medicines/download-medicine-data}

    \item Health Canada,
    \emph{Drug Product Database}.\\
    \url{https://www.canada.ca/en/health-canada/services/drugs-health-products/drug-products/drug-product-database.html}
\end{itemize}

\end{document}
